\documentclass[10pt,a4paper]{article}
% --- Codificación e idioma ---
\usepackage[utf8]{inputenc}
\usepackage[T1]{fontenc}
\usepackage[spanish,es-tabla]{babel}
% babel-spanish activa la particion silabica (shorthands "<> y otros). Dentro de
% entornos verbatim (lstlisting) eso rompe la compilacion. Se desactiva.
\usepackage{etoolbox}
\AtBeginDocument{\shorthandoff{<>}}
\hyphenpenalty=10000
\exhyphenpenalty=10000
% --- Layout y utilidades ---
\usepackage[a4paper,left=30mm,right=30mm,top=30mm,bottom=30mm]{geometry}
\usepackage{graphicx}
\usepackage{float}
\usepackage{booktabs}
\usepackage{array}
\usepackage{listings}
\usepackage{caption}
\usepackage{hyperref}
\usepackage{enumitem}
\usepackage{amsmath}
\usepackage[table]{xcolor}
%------------------------------------------------------------------------------
% ENCABEZADO Y PIE DE PÁGINA
%------------------------------------------------------------------------------
\usepackage{fancyhdr}
\usepackage{lastpage}
\pagestyle{fancy}
\setlength{\headheight}{27pt}
\renewcommand{\headrulewidth}{1pt}
\renewcommand{\footrulewidth}{1pt}
\rhead{\footnotesize \textsc{TP N\textsuperscript{o} 2\\ SDN: NAT con POX y Mininet}}
%\lhead{\includegraphics[width=2.5cm]{logo}}
\cfoot{\small \textsc{Página \thepage\ de \pageref*{LastPage}}}
\hypersetup{
    colorlinks=true,
    linkcolor=blue!50!black,
    urlcolor=blue!60!black,
    citecolor=blue!50!black,
}
% --- Estilo de código ---
\definecolor{codebg}{RGB}{248,248,248}
\definecolor{codecomment}{RGB}{106,153,85}
\definecolor{codekeyword}{RGB}{86,86,196}
\definecolor{codestring}{RGB}{163,101,38}
\lstdefinestyle{codestyle}{
    backgroundcolor=\color{codebg},
    commentstyle=\color{codecomment},
    keywordstyle=\color{codekeyword},
    stringstyle=\color{codestring},
    basicstyle=\ttfamily\footnotesize,
    breakatwhitespace=false,
    breaklines=true,
    captionpos=b,
    keepspaces=true,
    numbers=left,
    numberstyle=\tiny\color{gray},
    numbersep=6pt,
    showspaces=false,
    showstringspaces=false,
    showtabs=false,
    tabsize=4,
    frame=single,
    rulecolor=\color{gray!40},
    language=Python,
    morekeywords={self,None,True,False}
}
\lstset{style=codestyle}
%------------------------------------------------------------------------------
% CARÁTULA
%------------------------------------------------------------------------------
\renewcommand*{\maketitle}{
\thispagestyle{empty}
\begin{center}
%\includegraphics[width=11cm]{logo}
\vspace{1cm}
\textsc{\Large Redes (TA048)}\\
\textsc{\small Facultad de Ingeniería --- Universidad de Buenos Aires}\\[0.6cm]
{\huge \bfseries Trabajo Práctico N\textsuperscript{o} 2}\\[0.2cm]
{\Large \bfseries SDN: NAT con POX y Mininet}\\[0.3cm]
\rule{0.8\textwidth}{0.4pt}\\[0.4cm]
{\large NAT saliente con traducción por puertos (PAT)}\\[1 cm]
{\large 2\textsuperscript{do} cuatrimestre 2026}\\
{\large Comisión jueves}\\[1.5cm]
\begin{minipage}{0.7\textwidth}
\begin{flushleft}
\large
\textbf{Integrantes:}\\[0.5cm]
Alen Monti --- 108081\\[0.3cm]
Nadeska Millan --- 110779 \\[0.3cm]
Elías Cuba --- 106933\\[0.3cm]
Netanel Jamilis --- 99093 \\[0.3cm]
Sebastian Furnier --- 110129
\end{flushleft}
\end{minipage}
\vfill
{\large \today}
\end{center}
}
\begin{document}
\maketitle
\clearpage
% =====================================================================
\section{Introducción}
% =====================================================================
El presente trabajo práctico tiene como objetivo implementar un
\textbf{NAT saliente con traducción por puertos (PAT, \emph{Port Address
Translation})} en el contexto de \emph{Redes Definidas por Software} (SDN).

A diferencia de una arquitectura tradicional, en SDN el comportamiento de la
red no está distribuido en los dispositivos de forwarding, sino que es definido
por un \textbf{controlador} que se comunica con ellos mediante el protocolo
\textbf{OpenFlow}. Esto separa el \emph{plano de control} (decisiones de
ruteo/traducción, en el controlador) del \emph{plano de datos} (reenvío de
paquetes, en los switches).

La solución se desarrolla sobre el controlador \textbf{POX} y se valida con la
herramienta de emulación \textbf{Mininet}. El controlador provisto por la
cátedra (\texttt{protorouter.py}) implementaba solamente: procesamiento de
paquetes IPv4, detección de tráfico de la red privada, reescritura de MACs e
instalación de reglas de forwarding. Sobre esa base se extendió el código para
implementar:

\begin{itemize}[nosep]
    \item \textbf{Protocolo ARP dinámico}: respuesta de requests dirigidas al
          router, generación de requests para resolver MACs y mantenimiento de
          una tabla ARP con expiración.
    \item \textbf{NAT/PAT}: traducción de IP y puerto origen (TCP/UDP) de las
          conexiones iniciadas desde la red privada, con mantenimiento de
          estado y soporte para múltiples hosts y conexiones simultáneas que
          comparten una única IP pública.
    \item \textbf{Instalación de flujos}: por cada conexión se instalan reglas
          en ambos sentidos para que el switch realice la traducción y el
          forwarding sin intervención posterior del controlador, con timeouts
          para liberar recursos.
\end{itemize}

% =====================================================================
\clearpage
\section{Descripción general de la solución}
% =====================================================================
\subsection{Topología}
La topología extendida (\texttt{topo.py}) consta de un switch OpenFlow
(\texttt{s1}) con tres hosts conectados:

\begin{itemize}[nosep]
    \item \texttt{h1}: host en la red pública (servidor),
          \texttt{200.0.0.1/24}, gateway \texttt{200.0.0.254}.
    \item \texttt{h2}: host privado, \texttt{192.168.1.2/24}, gateway
          \texttt{192.168.1.254}.
    \item \texttt{h3}: segundo host privado, \texttt{192.168.1.3/24}, mismo
          gateway.
\end{itemize}

El controlador hace que \texttt{s1} actúe como un router L3 con dos
interfaces:

\begin{table}[H]
\centering
\begin{tabular}{lcc}
\toprule
\textbf{Interfaz} & \textbf{IP} & \textbf{MAC} \\
\midrule
Pública (puerto 1)   & \texttt{200.0.0.254}   & \texttt{00:00:00:aa:aa:aa} \\
Privada (puertos 2-3) & \texttt{192.168.1.254} & \texttt{00:00:00:bb:bb:bb} \\
\bottomrule
\end{tabular}
\caption{Interfaces del router implementado sobre el switch s1.}
\end{table}

Importante: \textbf{no existe ninguna entrada ARP estática} en ningún host;
además, ninguna MAC o IP de los hosts está hardcodeada en el controlador. La
dirección del servidor puede cambiarse (IP y/o MAC) y la solución sigue
funcionando, porque todo se descubre dinámicamente por ARP.

\subsection{Arquitectura del controlador}
El módulo \texttt{ext/protorouter.py} mantiene, por cada conexión de switch:

\begin{itemize}[nosep]
    \item una \textbf{tabla ARP dinámica} (\texttt{IP -> MAC}, con expiración a
          los 5 minutos);
    \item una \textbf{cola de resoluciones pendientes}: paquetes que no se
          pudieron reenviar porque falta la MAC destino y se enviarán cuando
          llegue la respuesta ARP;
    \item la \textbf{tabla NAT}, con un registro por conexión, indexada en
          ambos sentidos (clave saliente y clave entrante) para demuxear en
          $O(1)$;
    \item un \textbf{conjunto de puertos públicos en uso} para evitar
          colisiones;
    \item un \textbf{timer} periódico (10 s) que expira entradas ARP viejas y
          registros NAT cerrados.
\end{itemize}

El flujo de un paquete en el \texttt{PacketIn} es:

\begin{enumerate}[nosep]
    \item Si es \textbf{ARP}: se aprende de la tupla (IP, MAC) que traiga y se
          responde o inunda (\emph{flood}) según corresponda (sección~\ref{sec:arp}).
    \item Si es \textbf{IPv4}:
    \begin{itemize}[nosep]
        \item origen privado y destino público $\Rightarrow$ \emph{tráfico
              saliente}: se busca/crea el registro NAT, se traduce IP y puerto
              de origen, se resuelve la MAC destino por ARP si hace falta y se
              instala el par de flujos (sección~\ref{sec:nat}).
        \item destino = IP pública del router $\Rightarrow$ \emph{tráfico
              entrante}: se demuxea por la clave entrante de la tabla NAT, se
              traduce IP y puerto de destino hacia el host privado original y
              se reenvía por la interfaz privada.
        \item cualquier otra combinación se descarta.
    \end{itemize}
\end{enumerate}

% =====================================================================
\clearpage
\section{Manejo de ARP}
\label{sec:arp}
% =====================================================================
El protocolo ARP es el componente que permite al NAT resolver direcciones MAC
en ambas redes, sin entradas estáticas. La implementación tiene tres partes.

\subsection{Respuesta de requests dirigidas al router}
Cuando llega un ARP request cuyo \texttt{protodst} es una de las IPs del
router (\texttt{200.0.0.254} o \texttt{192.168.1.254}), el controlador responde
inmediatamente con un ARP \emph{reply} que usa la MAC de la interfaz
correspondiente (\texttt{aa:aa:aa...} para la pública, \texttt{bb:bb:bb...}
para la privada). Así, cualquier host que quiera comunicarse con ``su
gateway'' la resuelve dinámicamente.

\subsection{Generación de requests}
Cuando el NAT necesita reenviar un paquete y no conoce la MAC destino, envía
un ARP request desde la interfaz por donde saldría el paquete (con la IP y MAC
de esa interfaz como origen) y \textbf{enchufa} (\emph{queue}) el paquete. Al
llegar el reply, se aprende la MAC, se envían todos los paquetes encolados y
se completan las instalaciones de flujo que quedaron frenadas. Esto se
necesita tanto para el servidor público (\texttt{h1}) como para los hosts
privados (\texttt{h2}, \texttt{h3}).

\subsection{Tabla ARP dinámica y aprendizaje}
Todo paquete ARP (request o reply) que atraviesa el switch se usa para
\emph{aprender}: la tupla (\texttt{protosrc}, \texttt{hwsrc}) se guarda en la
tabla junto con el timestamp y el puerto del switch por el que se aprendió.
Las entradas expiran a los 5 minutos (\texttt{ARP\_ENTRY\_TIMEOUT}), y un
timer periódico hace la limpieza. Si llega una MAC distinta para una IP
conocida, se actualiza la entrada (y se lo indica en el log).

Los ARP requests dirigidos a otras IPs (por ejemplo, entre hosts del mismo
segmento) se inundan por el switch como en un aprendizaje convencional.

% =====================================================================
\clearpage
\section{Diseño de la tabla NAT (PAT)}
\label{sec:nat}
% =====================================================================
\subsection{Registro por conexión}
Cada conexión traducida se representa con la clase \texttt{NATRecord}:

\begin{table}[H]
\centering
\begin{tabular}{ll}
\toprule
\textbf{Campo} & \textbf{Significado} \\
\midrule
\texttt{proto}        & Protocolo: 6 (TCP) o 17 (UDP) \\
\texttt{private\_ip}  & IP del host interno \\
\texttt{private\_port}& Puerto origen del host interno \\
\texttt{remote\_ip}   & IP del host público (servidor) \\
\texttt{remote\_port} & Puerto destino en el servidor \\
\texttt{public\_port} & Puerto público asignado en el NAT \\
\bottomrule
\end{tabular}
\caption{Campos de un registro de la tabla NAT.}
\end{table}

La tabla se indexa en ambos sentidos:

\begin{itemize}[nosep]
    \item \textbf{clave saliente} $(proto,\ IP_{priv},\ port_{priv},\
          IP_{pub},\ port_{pub})$: para detectar si un paquete saliente
          pertenece a una conexión ya traducida;
    \item \textbf{clave entrante} $(proto,\ port_{público},\ IP_{pub},\
          port_{pub})$: para demuxear los paquetes que llegan del servidor y
          saber a qué host privado redirigirlos.
\end{itemize}

Además, cada registro guarda un timestamp y expira tras
\texttt{NAT\_RECORD\_TIMEOUT} (65 s, levemente mayor al \texttt{hard\_timeout}
de los flujos) para que no queden ``colgados'' registros de conexiones que ya
no tienen flujos en el switch.

\subsection{Ejemplo conceptual}
\begin{table}[H]
\centering
\begin{tabular}{lcccc}
\toprule
\textbf{Protocolo} & \textbf{IP privada} & \textbf{Puerto priv.} &
\textbf{IP pública} & \textbf{Puerto púb.} \\
\midrule
TCP & 192.168.1.2 & 42480 & 200.0.0.254 & 42480 \\
TCP & 192.168.1.3 & 50000 & 200.0.0.254 & 50000 \\
TCP & 192.168.1.2 & 42482 & 200.0.0.254 & 51000 \\
UDP & 192.168.1.3 & 5353  & 200.0.0.254 & 5353  \\
\bottomrule
\end{tabular}
\caption{Ejemplo de la tabla NAT con conexiones simultáneas de h2 y h3.}
\end{table}

Obsérvese que dos hosts distintos pueden incluso compartir el mismo puerto
público a la vez si van a servidores distintos: la clave entrante incluye la
IP del servidor, así que no hay ambigüedad.

\subsection{Asignación de puertos públicos}
La función \texttt{\_pick\_public\_port} implementa la política de asignación:

\begin{enumerate}[nosep]
    \item Si el puerto origen privado cae en el rango asignable
          $[10000, 65535]$ y no está en uso para ese protocolo,
          \textbf{se reutiliza} (evita traducciones innecesarias).
    \item Si no, se asigna un puerto del mismo rango elegido al azar,
          avanzando cíclicamente hasta encontrar uno libre. La elección
          aleatoria del punto de partida distribuye la carga sobre el espacio
          de puertos.
    \item El par $(proto, puerto)$ se marca como usado; se libera cuando el
          registro expira.
\end{enumerate}

Con esto se garantiza que dos conexiones simultáneas nunca obtengan el mismo
$(proto,\ puerto\_público)$ y que los recursos se reutilicen cuando las
conexiones terminan.

\subsection{Traducción de paquetes}
Para el tráfico saliente se reescriben: IP origen ($\rightarrow$
\texttt{PUBLIC\_IP}), puerto origen TCP/UDP ($\rightarrow$ puerto público
asignado) y MAC origen ($\rightarrow$ \texttt{PUBLIC\_MAC}); la MAC destino se
resuelve por ARP. Para el entrante, simétricamente: IP y puerto destino
($\rightarrow$ host privado original) y MAC destino del host privado.

Los checksums de IP y de la capa de transporte se recalculan al reempaquetar
el paquete con la librería de POX (\texttt{pack()}), que recompute el checksum
de IP y de TCP/UDP con los nuevos valores de los campos.

\subsection{ICMP}
El ping (ICMP echo) no tiene puertos, pero sí un \emph{identificador} de echo
que cumple la misma función de demultiplexado. El NAT lo trata de forma
análoga: usa el identificador como ``puerto'' para crear el registro, traduce
la IP origen y demuxea los echo replies entrantes por esa clave. Así,
\texttt{h2 ping h1} funciona y puede coexistir con los pings de \texttt{h3}.

% =====================================================================
\clearpage
\section{Estrategia de instalación de flujos}
\label{sec:flows}
% =====================================================================
El objetivo del diseño es que \textbf{el switch haga todo el trabajo} una vez
establecida la traducción, minimizando la intervención del controlador.

\subsection{Cuándo se instalan}
Los flujos de una conexión se instalan en el primer paquete de la conexión
(cuando se crea el registro), pero \textbf{solo cuando se conocen las dos
MACs} involucradas (la del servidor y la del host privado). Si alguna falta,
se dispara la resolución ARP y la instalación se completa al llegar la
respuesta. Esto evita instalar reglas con MACs incorrectas o vacías.

\subsection{Reglas instaladas}
Por cada conexión se envían dos \texttt{ofp\_flow\_mod} con \emph{match
exacto}:

\begin{itemize}[nosep]
    \item \textbf{Saliente}: \texttt{in\_port} = puerto privado donde entró el
          primer paquete, \texttt{dl\_type}=0x800, \texttt{nw\_src},
          \texttt{nw\_dst}, \texttt{nw\_proto}, \texttt{tp\_src} y
          \texttt{tp\_dst} de la conexión. Acciones:
          \texttt{SET\_NW\_SRC} (IP pública), \texttt{SET\_TP\_SRC} (puerto
          público, si cambia), \texttt{SET\_DL\_SRC/DST} (MACs) y
          \texttt{OUTPUT} por el puerto público.
    \item \textbf{Entrante}: \texttt{in\_port} = puerto público,
          \texttt{nw\_src/nw\_dst/tp\_src/tp\_dst} del flujo invertido visto
          desde el router. Acciones: \texttt{SET\_NW\_DST} (IP privada),
          \texttt{SET\_TP\_DST} (puerto privado, si cambia), reescritura de
          MACs y \texttt{OUTPUT} por el puerto privado del host.
\end{itemize}

Como el \emph{match} es exacto (incluye puertos), las reglas de distintas
conexiones no colisionan y la tabla del switch no crece con entradas
agregadas.

\subsection{Timeouts y liberación de recursos}
Cada flujo se instala con:
\begin{itemize}[nosep]
    \item \texttt{idle\_timeout = 10} s: si no pasa tráfico, la regla se elimina;
    \item \texttt{hard\_timeout = 60} s: límite máximo de vida de la regla;
    \item flag \texttt{OFPFF\_SEND\_FLOW\_REM}: el switch notifica al
          controlador la remoción, lo que permite (en futuras extensiones)
          liberar el puerto público en ese momento.
\end{itemize}
En paralelo, el registro NAT expira a los 65 s si no hay actividad, y el timer
periódico libera el puerto público. De esta forma, tanto las reglas del
switch como el estado del controlador se limpian solos.

\subsection{Tráfico que llega al controlador}
Con este diseño, el controlador solo ve: (a) paquetes ARP, (b) el primer
paquete de cada conexión, (c) paquetes cuya MAC destino todavía no se resolvió
y (d) tráfico cuya conexión ya expiró (que re-dispara la creación del
registro). Todo el tráfico de datos de una conexión establecida circula
únicamente por el switch.

% =====================================================================
\clearpage
\section{Decisiones de diseño}
% =====================================================================
\begin{enumerate}
    \item \textbf{Match exacto por conexión} en lugar de reglas agregadas por
          host: evita ambigüedades al demuxear el tráfico entrante y hace
          trivial la liberación de recursos por expiración. El costo es una
          entrada más en la tabla de flujos por conexión, aceptable para la
          escala del TP.
    \item \textbf{Doble indexación} de la tabla NAT (clave saliente y entrante
          apuntando al mismo registro): el demuxeo del tráfico entrante, que
          es la operación crítica, es $O(1)$ sin recorrer la tabla.
    \item \textbf{Cola de paquetes por resolución ARP}: en lugar de descartar
          el paquete que disparó el ARP request (lo que perdería el SYN de una
          conexión TCP, por ejemplo), se encola y se envía al llegar el reply.
          Esto hace que la primera conexión no pierda paquetes.
    \item \textbf{ICMP tratado como un ``flujo con puerto''}: el identificador
          de echo hace de puerto, lo que permite que los pings de varios hosts
          privados convivan sin colisiones, con el mismo mecanismo de PAT.
    \item \textbf{Reutilización del puerto origen} cuando es posible: reduce
          la cantidad de traducciones (solo se reescribe el puerto si cambió)
          y simplifica el debugging, porque el puerto visto en el servidor es
          el mismo que el del cliente.
    \item \textbf{No hardcodear nada de los hosts}: tanto las MACs de los
          hosts privados como la del servidor se descubren por ARP, y la IP
          del servidor se aprende del propio tráfico. Cambiar IP/MAC de
          cualquier host no rompe la solución (requisito de la demo).
\end{enumerate}

% =====================================================================
\clearpage
\section{Problemas encontrados y soluciones}
% =====================================================================
\begin{description}[style=nextline]
    \item[Respuestas que llegan antes que el flujo entrante.] Cuando el
    primer paquete saliente genera el ARP request hacia el servidor, el reply
    del servidor puede llegar antes de que el flujo entrante esté instalado.
    Solución: el registro NAT se crea \emph{antes} de instalar los flujos, así
    que aunque el paquete entre por \texttt{PacketIn} ya se puede demuxear; la
    instalación de flujos se completa al resolver la MAC.
    \item[Colisiones de puertos entre hosts.] Dos hosts privados pueden usar
    el mismo puerto origen a la vez. El conjunto de puertos en uso se lleva
    por $(proto,\ puerto)$, de modo que el segundo pedido recibe otro puerto
    automáticamente.
    \item[Checksums.] Cambiar IP o puerto invalida el checksum de IP y el de
    TCP/UDP. Se resuelve reempaquetando el paquete con la librería de POX,
    que recalcula ambos checksums con los valores finales.
    \item[ARP estático de la topología original.] La topología provista
    configuraba entradas ARP estáticas en los hosts para el debugging inicial.
    Se eliminaron por completo y se verificó que la primera comunicación de
    cada host resuelve su gateway por ARP dinámico.
    \item[Flujos con MACs desconocidas.] Instalar un flujo con la MAC destino
    aún sin resolver provocaba tráfico hacia MACs inválidas. Los flujos se
    instalan recién cuando se conocen ambas MACs; mientras tanto, el
    controlador procesa los paquetes de esa conexión.
\end{description}

% =====================================================================
\clearpage
\section{Pruebas realizadas}
% =====================================================================
\subsection{Conectividad básica}
Con el controlador y la topología corriendo:

\begin{lstlisting}[caption={Ping desde ambos hosts privados hacia el servidor.},captionpos=b]
mininet> h2 ping -c 3 200.0.0.1
mininet> h3 ping -c 3 200.0.0.1
\end{lstlisting}

Ambos pings funcionan. En el log del controlador se observa: (1) el ARP
request de cada host resolviendo su gateway, (2) la creación de los registros
ICMP en la tabla NAT y (3) la instalación de los flujos, tras lo cual los
siguientes paquetes del ping ya no generan \texttt{PacketIn}.

\subsection{Múltiples clientes}
En la ejecución registrada, \texttt{h2} y \texttt{h3} iniciaron conexiones TCP
hacia \texttt{200.0.0.1:9000}. El controlador asignó puertos públicos
distintos y demuxeó los paquetes de respuesta hacia la IP y el puerto privado
correspondientes. La evidencia disponible muestra ambos flujos a través del
NAT; no permite afirmar que las transferencias hayan sido simultáneas.

\subsection{Diferentes tipos de tráfico}
La salida adjunta registra tráfico ICMP echo (ida y vuelta) y TCP. No contiene
una prueba UDP, por lo que ese protocolo queda pendiente de verificación en
esta ejecución. Para TCP/UDP la traducción opera sobre los campos de puertos,
sin inspeccionar el payload.

\subsection{Resultados observados en la ejecución}
El log de POX confirma el aprendizaje ARP dinámico del servidor
\texttt{200.0.0.1} y de los hosts privados, seguido por la instalación de
flujos. Los mapeos NAT relevantes fueron:

\begin{table}[H]
\centering
\begin{tabular}{llrrl}
\toprule
\textbf{Protocolo} & \textbf{Host privado} & \textbf{Puerto/id privado} &
\textbf{Puerto/id público} & \textbf{Destino} \\
\midrule
ICMP & 192.168.1.2 & 3023 & 49381 & 200.0.0.1, secuencia 1 \\
ICMP & 192.168.1.3 & 3029 & 16179 & 200.0.0.1, secuencia 1 \\
TCP  & 192.168.1.2 & 54892 & 54892 & 200.0.0.1:9000 \\
TCP  & 192.168.1.3 & 39556 & 39556 & 200.0.0.1:9000 \\
\bottomrule
\end{tabular}
\caption{Traducciones observadas en el log de la ejecución compartida.}
\end{table}

Para ICMP se observan respuestas entrantes dirigidas de vuelta a
\texttt{192.168.1.2} y expiración posterior de los dos registros. En TCP, las
respuestas del servidor se reenviaron a \texttt{192.168.1.2:54892} y
\texttt{192.168.1.3:39556}, respectivamente. Los puertos TCP privados estaban
dentro del rango permitido y libres, por lo que se reutilizaron como puertos
públicos. El log de arranque también advierte que POX 0.7.0 declara soporte
para Python 3.6--3.9; la ejecución mostrada continuó con Python 3.12.3.

\subsection{Pruebas unitarias}
Se agregaron pruebas automatizadas para las rutas ICMP echo saliente y
entrante con dos hosts, la expiración y liberación de un registro NAT, el
aislamiento de los mapeos TCP y sus respuestas, y las entradas ARP de tres
campos (MAC, timestamp y puerto), incluyendo su expiración. Se ejecutan desde
la raíz del repositorio con:

\begin{lstlisting}[language=bash]
python3 -m unittest tests.unit.protorouter_test
\end{lstlisting}
La ejecución de este comando completó correctamente las tres pruebas.

\subsection{Independencia de valores hardcodeados}
Se modificaron la MAC y la IP del servidor (y de los hosts privados) en
\texttt{topo.py} y la solución siguió funcionando sin cambios en el
controlador, ya que todo el estado de MACs se aprende por ARP.

% =====================================================================
\section{Conclusiones}
% =====================================================================
El trabajo permite ver en concreto cómo una función clásica de red (el NAT)
puede implementarse como \emph{aplicación} sobre un controlador SDN: la
``inteligencia'' (crear traducciones, resolver ARP, decidir qué hacer con un
paquete nuevo) vive en el plano de control, y el trabajo repetitivo (traducir
y reenviar cada paquete de una conexión ya establecida) se delega al plano de
datos mediante reglas OpenFlow.

Los puntos que más aportan al aprendizaje son: el demuxeo del tráfico
entrante por $(puerto\_público,\ proto)$, que es exactamente lo que hace un
NAT real; la necesidad de mantener estado con expiración coherente entre el
controlador y los timeouts del switch; y la importancia de ARP como
infraestructura sobre la que se apoya todo el resto (sin resolución dinámica
de MACs, el NAT no puede reescribir tramas).

Una extensión natural sería liberar el puerto público al recibir el
\texttt{FlowRemoved} del switch (en lugar de por timer), y aplicar timeouts
distintos para TCP (según estado de la conexión) y UDP.

\end{document}
